\chapter{Adversarial Aftermath}
\label{short:adversarial-aftermath}

Never make victory perpetual.

When the contest ends, the common adversary may no longer hold the allies
together. What will their cooperation become when it must find another
reason to continue?

The notes ask whether the winners themselves must invoke the instrument of
defeat. That question remains unresolved. It asks us to examine what a
settlement demands of those who prevailed, as well as those who did not.

How can exceptional powers end? How can appeal, restoration, succession,
and negotiated exit remain possible? The aftermath belongs within the
design of cooperation, before victory tempts anyone to declare the design
finished.

\shortlimit{Cryptography does not make a settlement just or ensure that former allies remain united.}
\companionref{22}{chapter.22}{Governance notes; a settlement protocol remains to be developed.}
